\chapter{Geometria Kleina: przestrzeń i jej symetrie}
\label{ch:geometria-kleina}
\index{geometria Kleina}
\index{przestrzeń jednorodna}

Grupa Liego z Rozdziału~\ref{ch:grupy-liego} może działać na innej
rozmaitości. Zamiast pytać tylko o punkty i współrzędne, możemy
zapytać, które własności punktów, prostych lub odległości nie zmieniają
się pod tym działaniem. Taki był punkt wyjścia programu erlangeńskiego
Kleina: geometrię opisuje przestrzeń wraz z wybraną grupą jej
przekształceń. Ta sama rozmaitość może więc nieść różne geometrie.

W tym rozdziale zbudujemy język przestrzeni jednorodnych $G/H$.
Przykłady prowadzą od płaszczyzny euklidesowej i afinicznej przez
przestrzeń rzutową i sferę do torusa oraz dysku Poincarégo. Metryki
hiperboliczne wprowadziliśmy już w
Przykładzie~\ref{prz:trzy-modele-hiperboliczne}; tutaj obliczymy grupę
ich symetrii w wymiarze dwa.

\section{Działanie, orbita i stabilizator}
\label{sec:klein-dzialanie}

\begin{definicja}[Gładkie działanie grupy Liego]
\label{def:klein-dzialanie}
\index{działanie grupy Liego}
Niech $G$ będzie grupą Liego, a $M$ gładką rozmaitością bez brzegu.
\emph{Gładkim działaniem lewym} $G$ na $M$ nazywamy gładkie
odwzorowanie
\[
 a:G\times M\longrightarrow M,\qquad a(g,p)=g\cdot p,
\]
dla którego $e\cdot p=p$ oraz $(gh)\cdot p=g\cdot(h\cdot p)$.
Każdy $g$ daje dyfeomorfizm $a_g:p\mapsto g\cdot p$, którego
odwrotnością jest $a_{g^{-1}}$. Intuicyjnie grupa jest zbiorem
dozwolonych ruchów przestrzeni. Na przykład grupa $(\R^n,+)$ działa
na $\R^n$ przez translacje $b\cdot x=x+b$.
\end{definicja}

\begin{definicja}[Orbita, stabilizator i skuteczność]
\label{def:klein-orbita}
\index{orbita działania}
\index{stabilizator punktu}
\index{działanie skuteczne}
Dla $p\in M$ określamy \emph{orbitę} i \emph{stabilizator}
\[
 G\cdot p=\{g\cdot p:g\in G\},\qquad
 G_p=\{g\in G:g\cdot p=p\}.
\]
Działanie jest \emph{przechodnie}, gdy $G\cdot p=M$ dla pewnego,
a więc każdego $p\in M$. Jest \emph{skuteczne}, gdy jedynym elementem
ustalającym każdy punkt jest $e$. Orbita mówi, dokąd można przesunąć
punkt; stabilizator mówi, które ruchy pozostawiają go na miejscu.
Przy translacjach $\R^n$ na $\R^n$ orbita każdego punktu jest całą
przestrzenią, a stabilizator jest trywialny.
\end{definicja}

Stabilizator jest podgrupą: $e\in G_p$, a jeśli $g,h\in G_p$, to
$gh^{-1}\cdot p=g\cdot(h^{-1}\cdot p)=p$. Jest też domknięty,
gdyż jest przeciwobrazem punktu $p$ przez ciągłe odwzorowanie
$g\mapsto g\cdot p$, a rozmaitość $M$ jest Hausdorffa. Nie każda
orbita musi być domknięta ani być podrozmaitością z topologią
odziedziczoną z $M$; to rozróżnienie będzie potrzebne w twierdzeniu
poniżej.

\begin{lemat}[Stały rząd odwzorowania orbity]
\label{lem:klein-rzad-orbity}
Niech $p\in M$ i $o_p:G\to M$, $o_p(g)=g\cdot p$. Rząd
$\dd(o_p)_g$ nie zależy od $g$. Jego jądro w $e$ jest
$T_eG_p$, a $G_p$ jest zanurzoną podrozmaitością $G$ wymiaru
$\dim G-\operatorname{rank}\dd(o_p)_e$.
\end{lemat}
\begin{proof}
Dla $g,h\in G$ zachodzi $o_p\circ L_g=a_g\circ o_p$.
Różniczkując tę równość w $h$ i pamiętając, że $L_g$ oraz
$a_g$ są dyfeomorfizmami, dostajemy
\[
 \operatorname{rank}\dd(o_p)_{gh}
 =\operatorname{rank}\dd(o_p)_h.
\]
Odwzorowanie $o_p$ ma więc stały rząd $r$. Z
Twierdzenia~\ref{tw:rzad-staly} w mapach wokół $h\in G_p$ i $p$
ma ono postać $(u,v)\mapsto(u,0)$, gdzie $u\in\R^r$.
Jego przeciwobraz punktu $p$ jest lokalnie zadany równaniem $u=0$.
To dowodzi, że $G_p=o_p^{-1}(p)$ jest zanurzoną podrozmaitością
o podanym wymiarze i że $T_eG_p=\ker\dd(o_p)_e$.
Ograniczenia mnożenia i odwrotności w $G$ są gładkie, zatem
$G_p$ jest także grupą Liego z otrzymaną strukturą.
\end{proof}

Do opisu $G/H$ dla \emph{dowolnej} domkniętej podgrupy $H$ użyjemy
standardowego wyniku o grupach Liego. Jego założenie domkniętości
jest istotne: gęsta właściwa podgrupa może dać iloraz, w którym
nie da się rozdzielić dwóch klas otwartymi zbiorami.

\begin{twierdzenie}[Iloraz przez domkniętą podgrupę; wynik zewnętrzny]
\label{thm:klein-iloraz-liego}
Niech $G$ będzie skończenie wymiarową grupą Liego, a $H\subseteq G$
domkniętą podgrupą. Wtedy $H$ ma jednoznaczną strukturę zanurzonej
podgrupy Liego. Zbiór lewych warstw $G/H=\{gH:g\in G\}$ z topologią
ilorazową ma jednoznaczną strukturę gładkiej rozmaitości, dla której
rzutowanie $q:G\to G/H$ jest submersją. Zachodzi
\[
 \dim(G/H)=\dim G-\dim H,\qquad
 T_{eH}(G/H)\cong T_eG/T_eH.
\]
Działanie $G$ na $G/H$ przez $a\cdot(gH)=(ag)H$ jest gładkie
i przechodnie.
\end{twierdzenie}

Jest to połączenie twierdzenia o domkniętej podgrupie z twierdzeniem
o ilorazie przez wolne i właściwe działanie grupy Liego; dokładne
wersje i dowody podaje J.\ M.\ Lee, \emph{Introduction to Smooth
Manifolds}, rozdziały~20--21. W drugim twierdzeniu $H$ działa na
$G$ z prawej strony przez $(g,h)\mapsto gh$. Działanie jest wolne.
Aby sprawdzić jego właściwość, niech $C\subset G\times G$ będzie
zwarty i niech $K_1,K_2$ będą jego rzutami na czynniki.
Przeciwobraz $C$ przez $(g,h)\mapsto(g,gh)$ jest domknięty
w zwartym zbiorze
$K_1\times(H\cap K_1^{-1}K_2)$: drugi czynnik jest zwarty,
bo $K_1^{-1}K_2$ jest zwarty, a $H$ domknięta.
Zatem ten przeciwobraz jest zwarty. Wzór na wymiar i przestrzeń
styczną wynika już z lokalnej postaci submersji $q$ z
Wniosku~\ref{wn:submersja-lokalnie}, ponieważ
$q^{-1}(eH)=H$.

\begin{twierdzenie}[Orbita jako przestrzeń jednorodna]
\label{thm:klein-orbita-iloraz}
Niech grupa Liego $G$ działa gładko na rozmaitości $M$ bez brzegu
i niech $p\in M$, $H=G_p$. Wówczas
\[
 F:G/H\longrightarrow M,\qquad F(gH)=g\cdot p,
\]
jest dobrze określoną, różnowartościową immersją o obrazie
$G\cdot p$. Orbita ma przez $F$ kanoniczną strukturę gładkiej
rozmaitości z immersją w $M$, nawet jeśli jej topologia jako
podzbioru $M$ jest inna. Jeśli działanie jest przechodnie,
$F:G/H\to M$ jest
dyfeomorfizmem. Jeśli dana orbita jest zanurzoną podrozmaitością
$M$, to $F$ jest dyfeomorfizmem również na tę podrozmaitość.
\end{twierdzenie}
\begin{proof}
Z Lematu~\ref{lem:klein-rzad-orbity} stabilizator jest domkniętą
zanurzoną podgrupą, więc $G/H$ istnieje na mocy
Twierdzenia~\ref{thm:klein-iloraz-liego}. Równość
$gH=g'H$ oznacza $g'=gh$ dla pewnego $h\in H$; wtedy
$g'\cdot p=g\cdot(h\cdot p)=g\cdot p$. Odwrotnie, jeśli
$g'\cdot p=g\cdot p$, to $g^{-1}g'\in H$, zatem $g'H=gH$.
To daje poprawność definicji i różnowartościowość $F$.

Rzutowanie $q$ ma lokalne gładkie przekroje dzięki postaci
submersji. Ponieważ $F\circ q=o_p$, złożenie $o_p$ z takim
przekrojem daje gładkość $F$. Różniczkując w $e$, otrzymujemy
\[
 \dd F_{eH}\circ\dd q_e=\dd(o_p)_e.
\]
Jądra $\dd q_e$ i $\dd(o_p)_e$ są równe $T_eH$;
stąd $\dd F_{eH}$ jest różnowartościowe. Ta sama własność
w każdym $gH$ wynika z $G$-równoważności $F$ i z tego, że
działania przez $g$ na obu rozmaitościach są dyfeomorfizmami.
Zatem $F$ jest immersją.

Jeśli $G\cdot p=M$ i $\dim M>0$, stały rząd $o_p$ musi
równać się $\dim M$. Gdyby był mniejszy, każdy punkt $G$
byłby krytyczny, a Twierdzenie Sarda~\ref{thm:sard} mówiłoby,
że obraz $o_p(G)=M$ ma miarę zero w $M$. To niemożliwe,
bo niepusty otwarty zbiór w mapie ma dodatnią miarę.
Dla $\dim M=0$ wniosek o submersji jest natychmiastowy.
Zatem $\dd(o_p)_e$ jest surjektywne, a powyższa równość
różniczek pokazuje, że $\dd F_{eH}$, a więc każda
$\dd F_{gH}$, jest izomorfizmem. Twierdzenie o funkcji
odwrotnej czyni $F$ lokalnym dyfeomorfizmem. Jest ono także
bijekcją na $M$, więc jego lokalne odwrotności składają się
na gładką odwrotność globalną.

Jeśli zamiast $M$ rozważymy zanurzoną orbitę $N=G\cdot p$,
różniczki odwzorowania $F:G/H\to N$ mają ten sam rząd
$\dim(G/H)$, zaś argument Sarda zastosowany do przechodniego
działania na $N$ daje $\dim N=\dim(G/H)$. Stosujemy więc
poprzedni przypadek do $N$.
\end{proof}

Ostatnie założenie w twierdzeniu ma znaczenie. Działanie
$\R$ na $T^2$ przez
$t\cdot[x,y]=[x+t,y+\alpha t]$ dla niewymiernego $\alpha$
ma trywialny stabilizator, lecz orbita $[0,0]$ jest gęsta
i nie jest podrozmaitością z topologią odziedziczoną z torusa.
Jest to nawinięcie niewymierne opisane po
Przykładzie~\ref{prz:osemka}; odwzorowanie orbity jest
różnowartościową immersją, ale nie zanurzeniem w znaczeniu
Rozdziału~\ref{ch:rozmaitosci}.

\begin{definicja}[Geometria Kleina]
\label{def:klein-geometria}
\index{geometria Kleina!definicja}
\emph{Geometrią Kleina} będziemy nazywać parę $(G,X)$, w której
$X$ jest spójną gładką rozmaitością, a $G$ grupą Liego działającą
na niej gładko i przechodnio. Jej \emph{model skuteczny} ma ponadto
skuteczne działanie. Wybór $x_0\in X$ przedstawia przestrzeń jako
$X\cong G/G_{x_0}$; stabilizator zachowuje wybrany punkt.
Na przykład $(\R^n,\R^n)$ z translacjami jest geometrią Kleina,
w której stabilizator zera jest trywialny.
\end{definicja}

Niektórzy autorzy wymagają skuteczności już w definicji. Możemy
zawsze usunąć nieskuteczność: jądro działania jest
\begin{equation}\label{eq:klein-jadro}
 K=\{g:g\cdot x=x\text{ dla każdego }x\in X\}
   =\bigcap_{g\in G}gHg^{-1},\qquad H=G_{x_0}.
\end{equation}
Druga równość wynika z $x=g\cdot x_0$: element $a$ ustala
$x$ dokładnie wtedy, gdy $g^{-1}ag\in H$. Jest to domknięta
podgrupa normalna. Dlatego $G/K$ jest grupą Liego i działa
skutecznie oraz przechodnio na $X$. W szczególności stabilizator
$H$ nie musi być trywialny, aby całe działanie było skuteczne.

\section{Metryka a wybór grupy}
\label{sec:klein-metryka}

Struktura $G/H$ sama nie wybiera metryki. Można to uściślić,
korzystając z działania adjungowanego z
Definicji~\ref{def:lie-ad}.

\begin{stwierdzenie}[Kiedy istnieje metryka niezmiennicza]
\label{prop:klein-metryka-niezmiennicza}
Niech $H$ będzie domkniętą podgrupą $G$, a
$\mathfrak g=T_eG$, $\mathfrak h=T_eH$. Metryki riemannowskie
na $G/H$ niezmiennicze względem działania $G$ odpowiadają
dodatnio określonym iloczynom skalarnym na
$\mathfrak g/\mathfrak h$, niezmienniczym względem
indukowanego działania $\operatorname{Ad}(H)$.
\end{stwierdzenie}
\begin{proof}
Każda metryka jest wyznaczona przez swoje wartości w punktach.
Jeśli jest $G$-niezmiennicza, jej wartość w $eH$ musi być
niezmiennicza pod stabilizatorem $H$. Utożsamienie
$T_{eH}(G/H)\cong\mathfrak g/\mathfrak h$ z
Twierdzenia~\ref{thm:klein-iloraz-liego} przenosi działanie
$h\in H$ na $\operatorname{Ad}_h$ modulo $\mathfrak h$.
Jest ono dobrze określone, bo sprzężenie przez $h$
zachowuje $H$, więc $\operatorname{Ad}_h\mathfrak h=\mathfrak h$.
Istotnie,
\[
 h\exp_G(tX)H
 =h\exp_G(tX)h^{-1}H
 =\exp_G(t\operatorname{Ad}_hX)H.
\]
Różniczkowanie w $t=0$ daje tezę o reprezentacji na
przestrzeni stycznej.

Odwrotnie, dla takiego iloczynu skalarnego w $eH$
przenosimy go do $gH$ przez dyfeomorfizm działania $g$.
Jeśli $gH=g'H$ z $g'=gh$, dwa sposoby przeniesienia
różnią się działaniem $h$ w $eH$, więc dają ten sam wynik
właśnie dzięki niezmienniczości względem $H$.
Lokalne przekroje submersji $G\to G/H$ pokazują, że
powstała metryka zależy gładko od punktu.
\end{proof}

\section{Przestrzeń euklidesowa i afiniczna}
\label{sec:klein-euklidesowa-afiniczna}

\begin{przyklad}[Geometria euklidesowa]
\label{prz:klein-euklidesowa}
\index{grupa euklidesowa}
Niech $n\geq1$. Grupa ruchów euklidesowych
\[
 E(n)=\R^n\rtimes\mathrm O(n)
 =\{(b,A):b\in\R^n,\ A\in\mathrm O(n)\}
\]
ma iloczyn $(b,A)(c,B)=(b+Ac,AB)$ i działa na $\R^n$
według $(b,A)\cdot x=b+Ax$. Współrzędne $(b,A)$ oraz
gładkość działań w czynnikach czynią z $E(n)$ grupę Liego.
Translacje zapewniają przechodniość, a stabilizatorem $0$
jest $\{(0,A):A\in\mathrm O(n)\}\cong\mathrm O(n)$.
Twierdzenie~\ref{thm:klein-orbita-iloraz} daje
\[
 \R^n\cong E(n)/\mathrm O(n),\qquad
 \dim E(n)=n+\frac{n(n-1)}2,\qquad
 \dim\mathrm O(n)=\frac{n(n-1)}2.
\]
Działanie jest skuteczne: ruch ustalający $0$ ma $b=0$,
a ustalający następnie każdy wektor bazy ma $A=I$.
Macierz ortogonalna zachowuje $|x-y|$, więc zwykła metryka
euklidesowa jest $E(n)$-niezmiennicza. Są to wszystkie bijektywne
izometrie $\R^n$: po odjęciu obrazu zera izometria $f$ zachowuje
iloczyn skalarny z tożsamości polaryzacyjnej, a jej surjektywność
daje $\langle f(x+y)-f(x)-f(y),f(z)\rangle=0$ dla każdego $z$.
Jest więc addytywna, z ciągłości jednorodna, a także ortogonalna.
Dopuszczamy odbicia;
wybór $\R^n\rtimes\mathrm{SO}(n)$ daje wariant zachowujący
orientację.
\end{przyklad}

\begin{przyklad}[Grupa afiniczna i przestrzeń afiniczna]
\label{prz:klein-afiniczna}
\index{grupa afiniczna}
\index{przestrzeń afiniczna}
Niech
\[
 \operatorname{Aff}(n,\R)
 =\R^n\rtimes\mathrm{GL}(n,\R),\qquad
 (b,A)\cdot x=b+Ax.
\]
Wzór na iloczyn jest taki sam jak dla $E(n)$; odwrotność
wynosi $(b,A)^{-1}=(-A^{-1}b,A^{-1})$. Ponieważ
$\mathrm{GL}(n,\R)$ jest otwarta w przestrzeni macierzy,
te wzory są gładkie. Stabilizatorem $0$ jest
$\mathrm{GL}(n,\R)$, działanie jest przechodnie dzięki
translacjom i skuteczne z tego samego powodu co wyżej.
Zatem
\[
 \R^n\cong\operatorname{Aff}(n,\R)/\mathrm{GL}(n,\R),
 \qquad (n+n^2)-n^2=n.
\]
Jako rozmaitość jest to znowu $\R^n$, lecz większa grupa
zachowuje inne własności. Odwzorowanie $x\mapsto Ax+b$
wysyła proste na proste i zachowuje podział odcinka w danym
stosunku, gdyż
$(1-t)f(x)+tf(y)=f((1-t)x+ty)$. Nie musi zachować długości:
$x\mapsto2x$ należy do grupy afinicznej.
Właściwie żadna dodatnio określona metryka riemannowska
na $\R^n$ nie jest niezmiennicza względem całej tej grupy.
Jej stabilizator zera zawiera $x\mapsto2x$; niezmienniczość
w zerze wymagałaby $g_0(2v,2v)=g_0(v,v)$ dla każdego $v$,
co przeczy dodatniej określoności.
\end{przyklad}

\section{Przestrzeń rzutowa i sfera}
\label{sec:klein-rzutowa-sfera}

Przestrzeń $\mathbb{RP}^n$ z Przykładu~\ref{prz:RPn}
składa się z jednowymiarowych podprzestrzeni $\R^{n+1}$;
punkt zapisujemy $[x_0:\cdots:x_n]$. W tej geometrii istotne
są proste rzutowe i ich wzajemne położenie, a nie odległości.

\begin{przyklad}[Projektowa grupa liniowa]
\label{prz:klein-rzutowa}
\index{projektowa grupa liniowa}
Niech $n\geq1$ i
\[
 \mathrm{PGL}(n+1,\R)
 =\mathrm{GL}(n+1,\R)/(\R_*I),
 \qquad \R_*=\R\setminus\{0\}.
\]
Podgrupa skalarnych macierzy $\R_*I$ jest domknięta
w $\mathrm{GL}(n+1,\R)$, więc iloraz jest grupą Liego.
Klasa macierzy $A$ działa na $\mathbb{RP}^n$ przez
$[A]\cdot[x]=[Ax]$; zmiana $A$ przez niezerowy skalar
nie zmienia obrazu. W mapach rzutowych z
Przykładu~\ref{prz:RPn} działanie ma współrzędne będące
ilorazami funkcji liniowych o niezerowym mianowniku,
więc jest gładkie. Grupa $\mathrm{GL}(n+1,\R)$ przenosi
dowolną prostą liniową na dowolną inną, zatem działanie
jest przechodnie.

Wybierzmy $p=[1:0:\cdots:0]$. Jego stabilizator $P$
składa się z klas blokowych macierzy
\[
 \begin{pmatrix}\lambda&\beta\\0&B\end{pmatrix},
 \qquad \lambda\in\R_*,\quad
 \beta\in\R^{1\times n},\quad B\in\mathrm{GL}(n,\R).
\]
Warunek zerowego lewego dolnego bloku mówi dokładnie,
że pierwsza kolumna jest wielokrotnością $e_0$.
Stąd $\mathbb{RP}^n\cong\mathrm{PGL}(n+1,\R)/P$.
Kontrola wymiarów daje
\[
 \dim\mathrm{PGL}(n+1,\R)=(n+1)^2-1,
 \quad \dim P=n^2+n,
 \quad \dim\mathbb{RP}^n=n.
\]
Wymiar $P$ policzyliśmy jako
$1+n+n^2-1$: ostatnie $-1$ usuwa skalarne macierze.

Działanie $\mathrm{PGL}$ jest skuteczne. Jeśli $A$ zachowuje
każdą prostą, to $Ae_i=\lambda_i e_i$ dla wektorów bazy.
Zachowanie prostej $[e_i+e_j]$ wymusza
$\lambda_i=\lambda_j$, a więc $A$ jest skalarna i
$[A]$ jest jednością w $\mathrm{PGL}$.

Pełna geometria rzutowa nie ma niezmienniczej metryki
riemannowskiej. W mapie $[1:x_1:\cdots:x_n]$ element
$[\operatorname{diag}(1,2I_n)]\in P$ działa jako
$x\mapsto2x$. Różniczka w $p$ mnoży każdy wektor przez
$2$, co przeczy niezmienniczości dodatnio określonego
iloczynu skalarnego w $T_p\mathbb{RP}^n$. Można
\emph{wybrać} na $\mathbb{RP}^n$ metrykę pochodzącą ze
sfery, lecz jej grupa symetrii jest mniejsza niż pełne
$\mathrm{PGL}(n+1,\R)$.
\end{przyklad}

\begin{przyklad}[Sfera jako przestrzeń obrotów]
\label{prz:klein-sfera}
\index{sfera!przestrzeń jednorodna}
Dla $n\geq1$ grupa $\mathrm{SO}(n+1)$ działa na
$S^n\subset\R^{n+1}$ przez obroty. Działanie jest
przechodnie: jednostkowe wektory $x,y$ można uzupełnić
do dodatnio zorientowanych baz ortonormalnych, a macierz
przeprowadzająca jedną bazę w drugą leży w
$\mathrm{SO}(n+1)$. Stabilizator $e_0$ składa się z
$\operatorname{diag}(1,B)$, $B\in\mathrm{SO}(n)$.
Zatem
\[
 S^n\cong\mathrm{SO}(n+1)/\mathrm{SO}(n),
 \qquad
 \frac{n(n+1)}2-\frac{n(n-1)}2=n.
\]
Dla $n=1$ konwencja $\mathrm{SO}(1)=\{1\}$ sprawia,
że wzór mówi $S^1\cong\mathrm{SO}(2)$; nie ma tu
jednowymiarowego stabilizatora. Działanie jest skuteczne,
bo macierz ustalająca wszystkie punkty sfery ustala
każdy wektor przestrzeni. Metryka kulista otrzymana
przez ograniczenie metryki euklidesowej jest zachowana
przez $\mathrm{SO}(n+1)$.

Antypody $x$ i $-x$ na sferze wyznaczają tę samą
prostą liniową, stąd $S^n/\{x\sim-x\}\cong\mathbb{RP}^n$.
Metryka kulista schodzi na ten iloraz, ponieważ
antypodalność jest izometrią. Ta metryka jest dodatkową
strukturą: nie jest niezmiennicza względem pełnej grupy
rzutowej z poprzedniego przykładu. Grupa
$\mathrm O(n+1)$ działa na $\mathbb{RP}^n$ przez izometrie
tej metryki, z jądrem $\{I,-I\}$; jej skuteczny obraz
to $\mathrm{PO}(n+1)=\mathrm O(n+1)/\{I,-I\}$.
\end{przyklad}

\section{Torus: iloraz, grupa i geometria płaska}
\label{sec:klein-torus}

\begin{przyklad}[Dwa działania przechodnie na torusie]
\label{prz:klein-torus}
\index{torus!geometria Kleina}
Niech $n\geq1$ i $T^n=\R^n/\Z^n$, jak w
Przykładzie~\ref{prz:torus}. Rozdział~\ref{ch:grupy-liego}
wyposażył tę rozmaitość w strukturę grupy Liego:
$[x]+[y]=[x+y]$. Dlatego pierwszym, skutecznym modelem
jest działanie grupy $G=T^n$ na $X=T^n$ przez translacje.
Stabilizator $[0]$ to $H=\{[0]\}$ i
\[
 T^n\cong T^n/\{[0]\}.
\]

Możemy zamiast tego pozostawić grupę $G=(\R^n,+)$ i
zdefiniować $a\cdot[x]=[a+x]$. Działanie jest gładkie
i przechodnie, a stabilizator $[0]$ jest równy
$H=\Z^n$. Otrzymujemy zatem ten sam torus jako
\[
 T^n\cong\R^n/\Z^n.
\]
Tym razem działanie \emph{nie jest skuteczne}: każdy
$a\in\Z^n$ działa na każdym $[x]$ jak jedność.
Jądro~\eqref{eq:klein-jadro} wynosi właśnie $\Z^n$,
a skuteczna grupa ilorazowa to $\R^n/\Z^n=T^n$.
Podgrupa $\Z^n$ jest dyskretna, więc wzór na wymiar
ma postać $n-0=n$.

Zwykła metryka na $\R^n$ jest niezmiennicza względem
translacji przez $\Z^n$, toteż przechodzi na płaską
metrykę torusa. Metryka zależy od wybranego iloczynu
skalarnego na $\R^n$; samo przedstawienie torusa jako
grupy Liego nie wyznacza długości. Torus obrotowy
zanurzony w $\R^3$ z Przykładu~\ref{prz:torus-liego}
ma z kolei metrykę indukowaną przez to zanurzenie,
na ogół inną od płaskiej. Trzeba więc odróżniać
rozmaitość, jej strukturę grupową i wybraną metrykę.
\end{przyklad}

\section{Dysk Poincarégo jako iloraz grupy Liego}
\label{sec:klein-dysk}

\begin{przyklad}[Działanie $\mathrm{PSL}(2,\R)$]
\label{prz:klein-dysk}
\index{dysk Poincarégo}
Zacznijmy od górnej półpłaszczyzny
$\mathbb H^2=\{z=x+iy\in\C:y>0\}$. Niech
\[
 G=\mathrm{PSL}(2,\R)
 =\mathrm{SL}(2,\R)/\{I,-I\}.
\]
Podgrupa $\{I,-I\}$ jest domknięta i normalna,
więc $G$ jest grupą Liego. Jej wymiar wynosi trzy:
$\mathrm{SL}(2,\R)=\det^{-1}(1)$ jest regularną
hiperpowierzchnią w czterowymiarowej przestrzeni
macierzy, bo $\dd(\det)_A(A)=2\det A=2$;
iloraz przez dwuelementową podgrupę
nie zmienia wymiaru. Grupa $G$ działa na
$\mathbb H^2$ przez przekształcenia ułamkowo liniowe
\[
 \left[\begin{pmatrix}a&b\\c&d\end{pmatrix}\right]
 \cdot z=\frac{az+b}{cz+d},\qquad ad-bc=1.
\]
Znak całej macierzy nie zmienia ułamka. Mianownik
nie znika dla $z\in\mathbb H^2$, a bezpośrednie obliczenie daje
\begin{equation}\label{eq:klein-mobius-pochodna}
 \Im\frac{az+b}{cz+d}=\frac{y}{|cz+d|^2},
 \qquad
 \frac{\dd}{\dd z}\frac{az+b}{cz+d}
 =\frac1{(cz+d)^2}.
\end{equation}
Pierwsza równość pokazuje, że obraz pozostaje w
$\mathbb H^2$; druga, że działanie jest gładkie.

Działanie jest przechodnie. Dla $z=x+iy$ macierz
\[
 A_{x,y}=\begin{pmatrix}
 \sqrt y&x/\sqrt y\\0&1/\sqrt y
 \end{pmatrix}\in\mathrm{SL}(2,\R)
\]
wysyła $i$ w $x+iy$. Stabilizator $i$ obliczamy z
$(ai+b)/(ci+d)=i$: po porównaniu części rzeczywistej
i urojonej otrzymujemy $b=-c$ oraz $a=d$.
Warunek wyznacznika jest wtedy $a^2+c^2=1$.
Zatem stabilizator jest \emph{obrazem} $\mathrm{SO}(2)$
w $\mathrm{PSL}(2,\R)$, czyli
\[
 K=\mathrm{SO}(2)/\{I,-I\},\qquad
 \mathbb H^2\cong\mathrm{PSL}(2,\R)/K,
 \qquad 3-1=2.
\]
Grupa $K$ jest izomorficzna z kołem obrotów,
lecz nie jest dosłownie podgrupą macierzy
$\mathrm{SO}(2)\subset\mathrm{SL}(2,\R)$: w $G$
utożsamiliśmy macierze $I$ i $-I$. Konkretny izomorfizm
na zwykłą grupę obrotów ma postać $[R_\theta]\mapsto R_{2\theta}$.

Metryka hiperboliczna
$g_{\mathbb H}=((\dd x)^2+(\dd y)^2)/y^2$ z
Równania~\eqref{eq:hip-polprzestrzen} jest zachowana
przez $G$. Istotnie, przekształcenie holomorficzne
w~\eqref{eq:klein-mobius-pochodna} mnoży euklidesowy
element długości przez $|cz+d|^{-2}$, zaś wysokość
$y$ przez ten sam czynnik. W ilorazie definiującym
$g_{\mathbb H}$ czynniki te się skracają.
Jeśli element $G$ ustala wszystkie $z\in\mathbb H^2$,
to wielomian
$cz^2+(d-a)z-b$ znika na zbiorze otwartym, więc
$c=b=0$, $d=a$ i $a=\pm1$. W $G$ jest to jedność:
działanie jest skuteczne.

Przekształcenie Cayleya
\[
 C:\mathbb H^2\longrightarrow D,
 \qquad C(z)=\frac{z-i}{z+i},\qquad
 D=\{w\in\C:|w|<1\},
\]
jest dyfeomorfizmem na dysk, z odwrotnością
$C^{-1}(w)=i(1+w)/(1-w)$ i $C(i)=0$.
Z tożsamości
\[
 1-|C(z)|^2=\frac{4y}{|z+i|^2},
 \qquad |C'(z)|=\frac2{|z+i|^2}
\]
wynika przez podstawienie
\[
 C^*\!\left(\frac{4|\dd w|^2}{(1-|w|^2)^2}\right)
 =\frac{|\dd z|^2}{y^2}=g_{\mathbb H}.
\]
Tak otrzymujemy model dysku Poincarégo z
Równania~\eqref{eq:hip-kula} dla $n=2$ oraz jego
opis $D\cong\mathrm{PSL}(2,\R)/K$.
Po przeniesieniu działania przez $C$ stabilizator
zera działa obrotami dysku: dla macierzy obrotu $R_\theta$
bezpośrednie podstawienie daje
$C(R_\theta\cdot z)=e^{-2i\theta}C(z)$.
Brzeg $|w|=1$ nie należy
do przestrzeni; metryka sprawia, że droga ku niemu
ma nieskończoną długość.
\end{przyklad}

\section{Jak przechodzi się między geometriami}
\label{sec:klein-przejscia}

Grupy i ilorazy porządkują opisane przykłady, lecz
przejścia między nimi mają różny sens.

Po pierwsze, na \emph{tej samej} przestrzeni $\R^n$ mamy
inkluzję grup przekształceń
\[
 E(n)=\R^n\rtimes\mathrm O(n)
 \ \subset\ \operatorname{Aff}(n,\R)
 =\R^n\rtimes\mathrm{GL}(n,\R).
\]
Większa grupa oznacza mniej własności zachowywanych przez
wszystkie ruchy: długość jest euklidesowa, ale nie afiniczna;
współliniowość jest afiniczna. Oba modele mają tę samą
rozmaitość punktów, choć inne stabilizatory.

Po drugie, odwzorowanie $x\mapsto[1:x_1:\cdots:x_n]$
identyfikuje $\R^n$ z otwartym obszarem afinicznym
$U_0=\{[x_0:\cdots:x_n]:x_0\ne0\}$ w $\mathbb{RP}^n$.
Homomorfizm
\[
 (b,A)\longmapsto
 \left[\begin{pmatrix}1&0\\b&A\end{pmatrix}\right]
 \quad\text{z}\quad
 \operatorname{Aff}(n,\R)
 \longrightarrow\mathrm{PGL}(n+1,\R)
\]
jest różnowartościowy: równość dwóch klas wymagałaby
skalarnej proporcjonalności macierzy, a ich lewe górne
wyrazy równe $1$ wymuszają skalar $1$.
Realizuje zwykłe działanie
$x\mapsto Ax+b$ na $U_0$. Pełna grupa rzutowa
zawiera też przekształcenia, które nie zachowują
wybranego obszaru afinicznego: punkt z $\R^n$ może przejść
na hiperpłaszczyznę w nieskończoności. Zmienia się więc
zarówno grupa, jak i globalna przestrzeń.

Po trzecie, ilorazy zmieniają \emph{topologię} modelu.
Torus płaski powstaje z $\R^n$ przez identyfikację
$x\sim x+m$, $m\in\Z^n$, a przestrzeń rzutowa ze
sfery przez identyfikację punktów antypodalnych.
W obu przypadkach należy odróżnić grupę działającą na
przestrzeni pokrywającej od skutecznej grupy działającej
na ilorazie. Model dysku i półpłaszczyzny jest jeszcze
innym przejściem: przekształcenie Cayleya jest
izometrią między dwiema postaciami \emph{tej samej}
geometrii hiperbolicznej.

Geometrie trójwymiarowe omawiane przy twierdzeniu
Thurstona również można zapisywać jako przestrzenie
jednorodne grup Liego, lecz sama przechodniość działania
nie wystarcza do ich klasyfikacji: w tamtym zagadnieniu
ważne są także wybrana metryka i istnienie odpowiednich
ilorazów przez dyskretne grupy.

\par\smallskip
{\footnotesize\noindent\emph{Źródła i dalsza lektura:}
\href{https://math.ucr.edu/home/baez/erlangen/Klein-erlangen.pdf}
{F.\ Klein, \emph{A Comparative Review of Recent Researches in Geometry}}
(oryginalny program erlangeński z 1872 roku, przekład);
\href{https://link.springer.com/book/10.1007/978-1-4419-9982-5}
{J.\ M.\ Lee, \emph{Introduction to Smooth Manifolds}, wyd.~2},
rozdziały~20--21 (podgrupy domknięte, działania i ilorazy);
\href{https://u.math.biu.ac.il/~solomyb/TEACH/21/GrAct/Lec5b.pdf}
{B.\ Solomyak, \emph{The Hyperbolic Plane}},
wykład o działaniu $\mathrm{PSL}(2,\R)$ i stabilizatorze $i$.\par}
